\section{The Seven Enigmas of the World}

\begin{quotex}
\textit{Ignoramus et ignorabimus}

``we do not know and we shall never know"
\end{quotex}

In 1880, the physiologist \textbf{Emil du Bois-Reymond} proposed seven riddles of the world that science could not
answer and most likely never would. To wit:

\begin{enumerate}
\item the ultimate nature of matter and force 
\item the origin of motion 
\item the origin of life 
\item the purpose of living beings 
\item the origin of sensations 
\item the origin of thought and language 
\item free will 
\end{enumerate}
Of course, Bois-Reymond means a scientific explanation, that is, the discovery of laws based on observation of
phenomena, and not an a priori explanation from religious faith, metaphysics, or gnosis. Since it is difficult to
conceive what such a scientific explanation would even look like, scientists simply accept the ``big picture" of science
(e.g., evolution, quantum theory, ``big bang") and instead focus on technological issues. This is especially true in a
second tier class of intractable problems such as the origin and elimination of violence, inequality, racism, and other
socially destructive or undesirable aspects of human behaviour.

For example, the physicist Michio Kaku\footnote{\url{https://mkaku.org/}} asserted in a recent interview\footnote{\url{https://www.c-spanvideo.org/program/id/234213}} that a thousand years in the future, the human
race will have had sufficient time to solve and eliminate the conflicts resulting from religious, racial, economic or
national differences. Given Kaku's worldview, we have to assume that he means there is a
scientific explanation for such conflicts and, hence, a technological solution.

We are perhaps given a clue when he asserts that there is a gene responsible for generating mystical of supernatural
explanations of events. Unfortunately, there is absolutely no evidence for such a gene and his explanation of the
survival value of an allegedly incorrect explanation of events is less than convincing. But given that thesis, the
obvious conclusion to create a more rational (in Kaku's sense) humanity involves genetic
engineering or eugenics. If true, that would certainly solve the problem of religious conflict; does Kaku assume there
are specific genes that cause the other conflicts he mentions? If it is desirable to eliminate those genes or their
effects, given that the human race is genetically programmed, where does the impetus for that desire come from?

Well, early on, Kaku tells us that humans are ``naturally" scientific, so perhaps the science gene will gain the upper
hand. Unfortunately, Kaku inconsistently asserts that, unlike religious faith, doing science is actually not a natural
activity. It requires study and effort, qualities far from natural (or even common, presumably). Now we are coming
perilously close to the idea of science as transcendent knowledge. Then the physicist takes on the role of high priest,
asserting ``just so stories" of evolution and satisfying our mystical urges with strange stories of quantum parallel
universes.

So for what purpose? In order that humans, when no longer engaged in social conflicts, can focus on their real work of
exploring other planets. Again, why? Apparently, according to Kaku, for that same reason that a single virus will
infect a host organism and expand as rapidly and completely as it is able.

So even the physicist Kaku does not predict the solution to the seven enigmas in a thousand years. Now \textbf{Valentin Tomberg} states that Bois-Reymond should have said: ``given the current state of scientific knowledge and the scientific
method, these enigmas are insoluble." This leads us to these questions:

\begin{enumerate}
\item What do we accept as knowledge? 
\item What do we accept as method? 
\end{enumerate}
Science only accepts, as knowledge, the experience of the senses and the proposed laws that unify such experiences.
Tomberg instead proposes:

\begin{quotex}
you \emph{know} only what is verified by the concordance of all forms of experience in its totality
— experience of the senses, moral experience, mental experience, the collective experience of
other seekers for the truth, finally the experience of those whose knowledge has merited the title of wisdom and whose
will has been crowned with the title of sanctity. 

\end{quotex}
The method involves the observation of those experiences that the scientist ignores, viz., one's
thoughts, desires, emotions, sensations. The multiplicity of such phenomena is not resolved by a unitary law or
principle; rather it is necessary to become \emph{one} in oneself and \emph{one} with the spiritual world. Then a new
kind of experience — spiritual experience — will arise. And a new kind of
knowledge — cognition beyond discursive thought — will take place. Once there
is recognition of the relationship of all things and beings, then the corresponding method, the Method of Analogy, will
be understood.

This is the beginning of \textbf{Hermetic Science} and the enigmas will reveal their solutions.



\flrightit{Posted on 2010-10-09 by Cologero }

\begin{center}* * *\end{center}

\begin{footnotesize}\begin{sffamily}



\texttt{Andrew on 2010-10-11 at 02:33 said: }

More evidence that chemistry and the other sciences should never have divorced themselves from Alchemy.
Kaku's statements are a great illustration of the Fool Arcana, though.


\hfill

\texttt{James O'Meara on 2010-10-11 at 11:34 said: }

If you ask the Man in the Street why science is worthwhile, he'll mention atom bombs and [someday]
curing cancer. Unfortunately, it seems the ``best" scientists quickly tire of this practical work, feeling it to be
somewhat ``beneath" them, and contrive to reconfigure themselves as ``public figures" with important things to say about
metaphysics and morality. This must stop, and given the damage they do to society, I would be willing to toss out the
science [which is mostly ``zion-ce" anyway] as well. As Guenon said, the Western world neglects ``the one thing needful"
in favor of what other civilizations thought were contemptible toys.


\hfill

\texttt{Matt on 2010-10-11 at 13:33 said: }

Kaku is just one of many scientists who think that because they are a scientist it also means they are a philosopher,
and that science can just absorb the disciple of philosophy (or disregard it….like almost anything else that
doesn't fall under the usual sciences). I guess the one bright side to this is that it provides
those of us with some chuckles when these types make glaring philosophical and logical errors and are oblivious that
they made them.


\hfill

\texttt{kadambari on 2010-10-14 at 15:11 said: }

A philosophical education is different from mere science, as the physical sciences concern abstraction from phenomena,
they do not concern themselves with questions of living. Without any higher purpose placed upon it, science is
meaningless and directionless and clearly not an end in itself. All intelligent scientists realize this.Plato writes of
philosophy in the seventh letter:

One statement at any rate I can make in regard to all who have written or who may write with a claim to knowledge of the
subjects to which I devote myself–no matter how they pretend to have acquired it, whether from my instruction or from
others or by their own discovery. Such writers can in my opinion have no real acquaintance with the subject. I
certainly have composed no work in regard to it, nor shall I ever do so in the future, for there is no way of putting
it in words like other studies. Acquaintance must come rather after a long period of attendance on instruction in the
subject itself and of close companionship ( I suppose he means supervision of a correct ``teacher") when, suddenly, like
a blaze kindled by a leaping spark, it is generated in the soul and becomes self-sustaining.

Besides, this at any rate I know, that if there were to be a treatise or a lecture on this subject, I could do it best.
I am also sure for that matter that I should be very sorry so see such a treatise poorly written. If I thought it
possible to deal adequately with the subject in a treatise or a lecture for the general public, what finer achievement
would there have been in my life than to write a work of great benefit to mankind and to bring the nature of things to
light for all men? I do not, however, think the attempt to tell mankind of these matters a good thing, except in the
case of some few who are capable of discovering the truth for themselves with a little guidance. In the case of the
rest to do so would excite in some an unjustified contempt in a thoroughly offensive fashion, in others a certain lofty
and vain hopes as if they had acquired some awesome lore(341c-e)…

What is amazing is that he states that affinity for wisdom is something inborn and cannot be taught, which is what I am
beginning to think the older I get…

He goes on : To sum it all up in one word, natural intelligence and a good memory are equally powerless to aid the man
who has not an inborn affinity with the subject. Without such endowments there is of course not the slightest
possibility. Hence all who have no natural aptitude for and affinity with justice and all the other noble ideals,
though in the study of other matters they may be both intelligent and retentive–such will never any of them attain to
an understanding of the most complete truth in regard to moral concepts. The study of virtue and vice must be
accompanied by an inquirty into what is false and true of esixtence in general, and must be carried on my constant
practice throughout a long period….Hardly after practising detailed comparisons of names and definitions (I assume the
study of logic), and visual and other sense perceptions, after scrutinizing them in benevolent disputation by the use
of question and answer without jealously, at last in a flash understanding of each blazes up, and the mind as it exerts
all its powers to the limit of human capacity, is flooded with light…For this reason no serious man will ever think of
writing about serious realities for the general public so as to make them a prey to envy and perplexity. In a word, it
is an inevitable conclusion from this that when anyone sees anywehre the written work of anyone, whether that of a
lawgiver in his laws or whatever it may be in some other form, the subject treated cannot have been his most serious
concern–that is, if he is himself a serious man. His most serious interests have their abode somewhere in the noblest
region of the field of his activity….(344c-d)


\hfill

\texttt{johnson j on 2022-07-07 at 21:57 said: }

Everyone saying science could answer these riddles eventually looks totally foolish now even moreso than they did in
2010, because after 2020, anyone who doesn't know science can't even answer
questions about physical things has proven themselves to be braindead.


\end{sffamily}\end{footnotesize}
